\section{Experiments}
\label{sec:exp}



\begin{table}[t]
    \centering
    \caption{\textbf{Quantitative comparisons on DTU dataset under sparse view settings. }}
    \label{tab:dtu_results}
    \resizebox{1.0\linewidth}{!}{
    \begin{tabular}{l|ccc|ccc|ccc}
    \toprule
    \multirow{2}{*}{Methods} & \multicolumn{3}{c}{PSNR↑} & \multicolumn{3}{c}{SSIM↑} & \multicolumn{3}{c}{LPIPS↓} \\
  
    \cmidrule{2-10}
    & 3-view & 6-view & 9-view & 3-view & 6-view & 9-view & 3-view & 6-view & 9-view \\
    \midrule
    DietNeRF~\cite{jain2021dietnerf} & 11.85 & 20.63 & 23.83 & 0.633 & 0.778 & 0.823 & 0.214 & 0.201 & 0.173 \\
    RegNeRF~\cite{niemeyer2022regnerf} & 18.89 & 22.20 & 24.93 & 0.745 & 0.841 & 0.884 & 0.190 & 0.117 & \colorbox{thirdplace}{0.089} \\
    FreeNeRF~\cite{yang2023freenerf} & 19.52 & 23.25 & 25.38 & 0.787 & 0.844 & 0.888 & 0.173 & 0.131 & 0.102 \\
    SparseNeRF~\cite{wang2023sparsenerf} & 19.47 & - & - & 0.829 & - & - & 0.183 & - & - \\
    % ReconFusion~\cite{wu2024reconfusion} & 20.74 & 23.62 & 24.62 & 0.875 & 0.904 & 0.921 & 0.124 & 0.105 & 0.094 \\
    % MuRF~\cite{xu2024murf} & 21.31 & 23.74 & 25.28 & 0.885 & 0.921 & 0.936 & 0.127 & 0.095 & 0.084 \\
    \midrule
    3DGS~\cite{kerbl20233d} & 10.99 & 20.33 & 22.90 & 0.585 & 0.776 & 0.816 & 0.313 & 0.223& 0.173 \\
    FSGS~\cite{zhu2024fsgs} & 17.34 & 21.55 & 24.33 & {0.818} & \colorbox{thirdplace}{0.880} & \colorbox{thirdplace}{0.911} & 0.169 & 0.127 & 0.106 \\
    DNGaussian~\cite{li2024dngaussian} & 18.91 & 22.10 & 23.94 & 0.790 & 0.851 & 0.887 & 0.176 & 0.148 & 0.131 \\

    CoR-GS~\cite{zhang2024cor} & {19.21} & \colorbox{secondplace}{24.51} & \colorbox{secondplace}{27.18} & 0.853 & \colorbox{secondplace}{0.917} & \colorbox{secondplace}{0.947} & 0.119 & \colorbox{secondplace}{0.068} & \colorbox{firstplace}{0.045} \\

    BinocularGS~\cite{han2024binocular} & \colorbox{secondplace}{20.71} & \colorbox{thirdplace}{24.31} & \colorbox{thirdplace}{26.70} & \colorbox{thirdplace}{0.862} & \colorbox{secondplace}{0.917} & \colorbox{secondplace}{0.947} & \colorbox{thirdplace}{0.111}& \colorbox{thirdplace}{0.073} & \colorbox{secondplace}{0.052} \\
  
  
    % ComapGS~\cite{} & 19.34 & 0.571 & 0.384 & 20.91 & 0.647 & 0.308 & 22.35 & 0.715 & 0.256 \\
    NexusGS~\cite{zheng2025nexusgs} & \colorbox{thirdplace}{20.21} & - & - & \colorbox{secondplace}{0.869} & - & - & \colorbox{secondplace}{0.102} & - & - \\
    % DropGaussians~\cite{} & 19.34 & 0.571 & 0.384 & 20.91 & 0.647 & 0.308 & 22.35 & 0.715 & 0.256 \\
    
  
    \midrule
    \textbf{Ours} & \colorbox{firstplace}{\textbf{21.77}} & \colorbox{firstplace}{\textbf{25.09}} & \colorbox{firstplace}{\textbf{27.19}} & \colorbox{firstplace}{\textbf{0.888}} & \colorbox{firstplace}{\textbf{0.928}} & \colorbox{firstplace}{\textbf{0.953}} & \colorbox{firstplace}{\textbf{0.092}} & \colorbox{firstplace}{\textbf{0.064}} & \colorbox{firstplace}{\textbf{0.045}} \\
    \bottomrule
    \end{tabular}
    }
    \vspace{-0.3cm}
  \end{table}
  
  
  \begin{figure*}[t]
    % \vspace{-0.3cm}
    \begin{center}
       \includegraphics[trim={0cm 0cm 0cm 0cm},clip,width=0.78\linewidth]{figures/DTU_com.pdf}
    \end{center}
    \vspace{-0.6cm}
    % \caption{ \textbf{Motivation.} }
    \caption{\textbf{Visual comparison on DTU~\cite{jensen2014large} dataset.}}
    % \label{fig:motivation}
    \label{fig:dtucom}
    \vspace{-0.3cm}
    \end{figure*}





\subsection{Datasets}
We conduct experiments on three standard benchmarks: LLFF~\cite{mildenhall2019local} 
for forward-facing scenes, DTU~\cite{jensen2014large}, a challenging dataset with 
extensive weakly-textured regions, for object-centric captures,  
and Blender~\cite{mildenhall2021nerf} for 360° object-centric scenes.
Following~\cite{han2024binocular,zhang2024cor,jang2025comapgs}, 
we use 3, 6, 9 training views for LLFF/DTU and 8 views for Blender scenes.
Input images are downsampled by 8× (LLFF), 4× (DTU), and 2× (Blender) to balance 
quality and efficiency, consistent with prior work~\cite{han2024binocular, zhang2024cor,jang2025comapgs}.


% \subsection{Implementation Details}
% We build upon BinocularGS~\cite{han2024binocular} with dense point cloud 
% initialization for LLFF and DTU, and random initialization for Blender.
% We train for 30k iterations on LLFF and DTU, and 7k iterations on Blender, 
% consistent with the baseline.
% Geometric regularization begins at iteration 20k for LLFF/DTU and 4k for 
% Blender, while geometry-guided appearance optimization starts at iteration 
% 25k and 5k respectively. Both operate at every iteration once activated.
% All experiments run on an NVIDIA L40s GPU.
% Loss weights are set to $\lambda_{\text{mpc}} = 0.1$, 
% $\lambda_{\text{smooth}} = 0.01$, and $\lambda_{\text{app}} = 1.0$.
% All warping operations are accelerated using batched parallel processing in PyTorch.
% We report average results over three independent runs.

\subsection{Implementation Details}
We build upon the BinocularGS framework~\cite{han2024binocular} with dense point cloud 
initialization for LLFF and DTU, and random initialization for Blender.
We train for 30k iterations on LLFF and DTU datasets, and 7k iterations on Blender, 
consistent with the baseline.
Geometric regularization begins at iteration 20k for LLFF/DTU and 4k for 
Blender, while geometry-guided appearance optimization starts at iteration 
25k and 5k respectively. Both operate at every iteration once activated.
All experiments run on the NVIDIA L40s GPU.
Loss weights are set to $\lambda_{\text{mpc}} = 0.1$, 
$\lambda_{\text{smooth}} = 0.01$, and $\lambda_{\text{app}} = 1.0$.
All warping operations are accelerated using batched parallel processing in PyTorch.
We report average results over three independent runs with different seeds.


\begin{table}[t]
    \centering
    \caption{\textbf{Quantitative comparison on Blender~\cite{mildenhall2021nerf} for 8 views.}}
    \label{tab:blender}
    \resizebox{0.60 \columnwidth}{!}{
    \begin{tabular}{l ccc}
    \toprule
    \textbf{Methods} & \textbf{PSNR$\uparrow$} & \textbf{SSIM$\uparrow$} & \textbf{LPIPS$\downarrow$} \\
    \midrule
    DietNeRF~\cite{jain2021dietnerf} & 22.50 & 0.823 & 0.124 \\
    RegNeRF~\cite{niemeyer2022regnerf} & 23.86 & 0.852 & 0.105 \\
    FreeNeRF~\cite{yang2023freenerf} & 24.26 & 0.883 & 0.098 \\
    SparseNeRF~\cite{wang2023sparsenerf} & 22.41 & 0.861 & 0.199 \\
    \midrule
    3DGS~\cite{kerbl20233d} & 22.23 & 0.858 & 0.114 \\
    FSGS~\cite{zhu2024fsgs} & 22.76 & 0.829 & 0.157 \\
    DNGaussian~\cite{li2024dngaussian} & 24.31 & 0.886 & \colorbox{secondplace}{0.088} \\
    CoR-GS~\cite{zhang2024cor} & 23.98 & \colorbox{secondplace}{0.891} & 0.094 \\
    BinocularGS~\cite{han2024binocular} & \colorbox{thirdplace}{24.71} & 0.872 & 0.101 \\
    % ComapGS~\cite{} & 23.67 & 0.853 & 0.102 \\
    NexusGS~\cite{zheng2025nexusgs} & 24.37 & \colorbox{firstplace}{0.893 }& \colorbox{firstplace}{0.087} \\
    DropGaussians~\cite{park2025dropgaussian} &\colorbox{secondplace}{25.42} & \colorbox{thirdplace}{0.888} & \colorbox{thirdplace}{0.089} \\
    \midrule
    \textbf{Ours} & \colorbox{firstplace}{\textbf{25.56}} & {\textbf{0.884}} & {\textbf{0.100}} \\
    \bottomrule
    \end{tabular}
    }
    \vspace{-0.3cm}
    \end{table}
  
    \begin{figure}[t]
      \vspace{-0.3cm}
      \begin{center}
         \includegraphics[trim={0cm 0cm 0cm 0cm},clip,width=0.8\linewidth]{figures/Blender_com.pdf}
      \end{center}
      \vspace{-0.3cm}
      % \caption{ \textbf{Motivation.} }
      \caption{\textbf{Visual comparison on Blender~\cite{mildenhall2021nerf} dataset.}}
      % \label{fig:motivation}
      \label{fig:blendercom}
      \vspace{-0.3cm}
      \end{figure}


\subsection{Baselines}
We compare against state-of-the-art methods from both NeRF and 3D Gaussian 
Splatting (3DGS) paradigms.
For NeRF-based methods, we include DietNeRF~\cite{jain2021dietnerf}, 
RegNeRF~\cite{niemeyer2022regnerf}, FreeNeRF~\cite{yang2023freenerf}, 
and SparseNeRF~\cite{wang2023sparsenerf}.
For 3DGS-based methods, we evaluate vanilla 3DGS~\cite{kerbl20233d} and 
sparse-view variants: FSGS~\cite{zhu2024fsgs}, DNGaussian~\cite{li2024dngaussian}, 
CoR-GS~\cite{zhang2024cor}, BinocularGS~\cite{han2024binocular}, 
DropGaussians~\cite{park2025dropgaussian}, NexusGS~\cite{zheng2025nexusgs}, 
and ComapGS~\cite{jang2025comapgs}.
% Quantitative results are taken from the original papers.
% , while qualitative 
% comparisons are based on our reproductions.

\subsection{Comparisons}
\noindent \textbf{Results on LLFF~\cite{mildenhall2019local} .}
\cref{tab:llff_results} presents quantitative results on the LLFF dataset. 
Our method achieves state-of-the-art performance across all view settings: 
\textbf{+0.76 dB} improvement at 3 views, \textbf{+0.17 dB} over 
ComapGS~\cite{jang2025comapgs} at 6 views, and competitive performance 
at 9 views (\textbf{+0.24 dB} over the third-best). \cref{fig:llffcom} 
demonstrates fewer artifacts in fine details and sharper object boundaries. 
Notably, our depth maps exhibit significantly clearer geometric edges in 
both background textures and foreground structures.

\noindent \textbf{Results on DTU~\cite{jensen2014large}.} 
\cref{tab:dtu_results} shows consistent improvements across all settings: 
\textbf{+1.06 dB}, \textbf{+0.58 dB}, and \textbf{+0.01 dB} over baselines 
at 3, 6, and 9 views, respectively. \cref{fig:dtucom} shows our method 
renders clearer textures in regions (red boxes) via occlusion-aware 
photometric consistency (\cref{subsec:geometry}), while depth maps reveal 
sharper boundaries and finer details (white boxes).


\noindent \textbf{Results on Blender~\cite{mildenhall2021nerf}.} 
For 360° object-centric evaluation, \cref{tab:blender} reports results 
on the Blender dataset~\cite{mildenhall2021nerf} under 8-view setting. 
Our method achieves the best PSNR. While SSIM and LPIPS are slightly lower 
than some methods, this trade-off reflects our prioritization of geometric 
accuracy over perceptual optimization. \cref{fig:blendercom} shows 
fine-grained textures and geometrically accurate structures, 
validating that our geometric regularization (\cref{subsec:geometry}) 
and geometry-guided appearance optimization (\cref{subsec:appearance}) 
effectively preserve structural fidelity critical for 3D reconstruction.






\subsection{Analysis}

\noindent \textbf{Ablation Study.}  
% \subsection{Ablation Study}n
\cref{tab:ablation} validates each component's contribution on LLFF (3 views) and DTU (3 views) by removing: robust multi-view photometric 
consistency Loss \textbf{$\mathcal{L}_{\text{mpc}}^{\text{Fea}}$} (introduced in \cref{eq:mpc_final}), depth smoothness Loss \textbf{$\mathcal{L}_{\text{smooth}} $} (introduced in \cref{eq:smooth}), cycle consistency depth filtering (\textbf{CCDF}, introduced in \cref{eq:depth_mask}), 
and virtual-view photometric consistency Loss \textbf{$\mathcal{L}_{\text{app}}$} (introduced in \cref{eq:aug_loss}).All experiments use BinocularGS~\cite{han2024binocular} as the baseline model.

Removing \textbf{$\mathcal{L}_{\text{mpc}}^{\text{Fea}}$} causes severe degradation 
(LLFF: -0.38 dB, DTU: -0.46 dB). 
As shown in \cref{fig:ablation}, this leads to noticeable blur and noise in both 
RGB and depth, confirming that robust geometric regularization is essential.
\textbf{$\mathcal{L}_{\text{smooth}}$} contributes -0.10 dB on DTU and prevents 
loss of structural information in depth maps (\cref{fig:ablation}), effectively 
mitigating artifacts from depth discontinuities.
\textbf{CCDF} improves virtual view synthesis (DTU: -0.52 dB, 
LLFF: -0.34 dB without it). 
Without CCDF, visible rendering artifacts appear (\cref{fig:ablation}), indicating 
insufficient appearance-geometry consistency.
\textbf{$\mathcal{L}_{\text{app}}$} contributes -0.57 dB on DTU and maintains image 
sharpness (\cref{fig:ablation}), demonstrating that appearance optimization benefits 
from explicit geometric guidance.
Our full model integrates all components to achieve superior geometric fidelity 
with clear structures and crisp boundaries.


\noindent \textbf{Limitations.}  
% \noindent\textbf{View-dependent effects.}
Our method assumes view-independent appearance during virtual view synthesis. 
In regions with strong view-dependent effects (e.g., specular highlights and 
reflections), the warped appearance may provide incorrect supervision.
However, due to view sparsity, prior methods also struggle in these regions. 
We refer readers to the supplementary material for visual examples.
% \noindent\textbf{Computational cost.}
Despite our accelerated implementation, it requires additional 
computation, resulting in a training time approximately 1.5× that of the baseline. 
The memory overhead is modest (~0.3 GB on 3-view LLFF scenes).
We believe these costs are acceptable given the improvements in 
rendering quality and geometric accuracy under sparse views.

\begin{table}[t]
    \centering
    \caption{\textbf{Ablation studies on LLFF~\cite{mildenhall2019local} and DTU~\cite{jensen2014large}.}}
    \label{tab:ablation}
    \resizebox{1.0\linewidth}{!}{
    % \begin{tabular}{cccc|ccc|ccc|ccc}
      \begin{tabular}{cccc|ccc|ccc}
    \toprule
    \multirow{2}{*}{
    $\mathcal{L}_{\text{mpc}}^{\text{Fea}}$
    } & \multirow{2}{*}{
    $\mathcal{L}_{\text{smooth}} $
    } & \multirow{2}{*}{\textbf{CCDF}} & \multirow{2}{*}{
    $\mathcal{L}_{\text{app}}$
    } & \multicolumn{3}{c|}{\textbf{LLFF}(3-view)} & \multicolumn{3}{c}{\textbf{DTU}(3-view)}
    %  & \multicolumn{3}{c}{\textbf{Blender}(8-view)}  
     \\

    \cmidrule(lr){5-7} \cmidrule(lr){8-10} 
    % \cmidrule(lr){11-13}
    & & & & \textbf{PSNR}$\uparrow$ & \textbf{SSIM}$\uparrow$ & \textbf{LPIPS}$\downarrow$ & \textbf{PSNR}$\uparrow$ & \textbf{SSIM}$\uparrow$ & \textbf{LPIPS}$\downarrow$ 
    % & \textbf{PSNR}$\uparrow$ & \textbf{SSIM}$\uparrow$ & \textbf{LPIPS}$\downarrow$ 
    \\
    \midrule
   $\times$       &$\times$       & $\times$      & $\times$      & {21.44} & {0.751} & {0.168} & {20.71} & {0.862}& {0.111}
    %  & {24.71} & {0.872} & {0.101} 
     \\

    % \midrule
  %   \checkmark 
  % & & & & 21.44 & 0.751 & 0.168 & 20.71 & 0.862& 0.111 & 24.71 & 0.872 & 0.101 \\
 $\times$       &\checkmark&\checkmark &\checkmark & 21.82 & 0.761 & 0.165 & 21.31 & 0.883 & 0.099 
  % & 25.29 & 0.883 & 0.098 
  \\
  \checkmark& $\times$       &\checkmark &\checkmark& 22.16 & 0.775 & 0.159 & 21.67 & 0.879 & 0.093
  %  & 25.38 & 0.886 & 0.100 
   \\

  \checkmark  &  \checkmark  &  $\times$       & \checkmark  & 21.86 & 0.767 & 0.162 & 21.25 & 0.875& 0.104
  %  & 24.93 & 0.882 & 0.095
   \\

  \checkmark  & \checkmark  & \checkmark & $\times$       & 21.79 & 0.763 & 0.164 & 21.20 & 0.870& 0.105
  %  & 24.82 & 0.878 & 0.096 
   \\

    \midrule
  \checkmark&\checkmark&\checkmark &\checkmark & {\textbf{22.20}} & {\textbf{0.778}} & {\textbf{0.157}} & {\textbf{21.77}} & {\textbf{0.888}} & {\textbf{0.092}} 
  % & {\textbf{25.56}} & {\textbf{0.894}} & {\textbf{0.089}} 
  \\
    \bottomrule
    \end{tabular}
    }
    % \vspace{-0.4cm}
  \end{table}  


  \begin{figure}[t]
    \vspace{-0.3cm}
    \begin{center}
       \includegraphics[trim={1cm 1cm 1cm 0cm},clip,width=1.0\linewidth]{figures/ablation.pdf}
    \end{center}
    \vspace{-0.8cm}
    % \caption{ \textbf{Motivation.} }
    \caption{\textbf{Visualization of ablation study results using 3-views.}}
    % \label{fig:motivation}
    \label{fig:ablation}
    \vspace{-0.3cm}
    \end{figure}
